\section{Controller Design}\label{sec:controller}

This section separates the observer result from the delayed-feedback result. The
latter is proved exactly; the observer is required to satisfy the explicit
finite-time interface stated in Lemma~\ref{lem:observer_interface}.

\subsection{Disturbance-observer interface}\label{sec:ptdo}
Let $\chi=\dot q_m$ and define
\begin{equation}
 \dot\chi=u_o+\Delta_a,\qquad
 u_o=M_e^{-1}(q_m)[\tau-C_e(q_m,\dot q_m)\dot q_m],\qquad
 \Delta_a=M_e^{-1}(q_m)d.
 \label{eq:observer_channel}
\end{equation}
The following assumption is the precise interface needed by the controller.
\begin{assumption}[Observer interface]\label{as:observer_interface}
There are $T_o>0$ and an estimate $\hat\Delta_a$ such that, under the
continuous-time observer, $\hat\Delta_a(t)=\Delta_a(t)$ for all $t\geq T_o$.
Before $T_o$, the observation error is bounded. The estimate is formed as
$\hat d=M_e\hat\Delta_a$.
\end{assumption}

For completeness, the observer used in the simulations is the prescribed-time
homogeneous observer
\begin{equation}
\begin{aligned}
 \dot z_1&=z_2+u_o+\ell_1(t)\,\phi_1(\varepsilon_1/\sigma),\\
 \dot z_2&=\ell_2(t)\,\phi_2(\varepsilon_1/\sigma),
 \qquad \varepsilon_1=\chi-z_1,
\end{aligned}\label{eq:ptdo}
\end{equation}
where
\begin{align}
 \phi_1(x)&=\operatorname{sig}^{1-\eta/2}(x)+
             \operatorname{sig}^{1+\eta/2}(x),\\
 \phi_2(x)&=\operatorname{sig}^{2-\eta}(x)+
             \operatorname{sig}^{2+\eta}(x)+\sigma\operatorname{sgn}(x),
 \end{align}
$0<\eta<1$, and $\ell_1,\ell_2$ are selected according to the
prescribed-time homogeneous-observer design in~\cite{jiangPrescribedtimeDisturbanceObserver2024}.
The gains must be selected using the disturbance regularity bounds required by
that design; in particular, Assumption~\ref{as:disturbance} alone is not
sufficient to establish exact convergence of~\eqref{eq:ptdo}.

\begin{lemma}[Observer interface lemma]\label{lem:observer_interface}
Suppose the hypotheses of the prescribed-time homogeneous-observer theorem
used to select $\ell_1,\ell_2$ hold, including the required bound on
$\dot\Delta_a$. Then the Filippov solution of~\eqref{eq:ptdo} satisfies
$z_2(t)=\Delta_a(t)$ for all $t\geq T_o$.
\end{lemma}
\begin{proof}
The error variables $\varepsilon_1=\chi-z_1$ and
$\varepsilon_2=\Delta_a-z_2$ satisfy the homogeneous observer error system
obtained by subtracting~\eqref{eq:ptdo} from~\eqref{eq:observer_channel}.
The cited prescribed-time homogeneous-observer theorem applies to this
system because its vector field is homogeneous in the two error variables and
its perturbation is $\dot\Delta_a$. It provides a Lyapunov function $V$ and
positive constants $c_1,c_2$ for which the time-rescaled error system obeys
$\dot V\leq-c_1V^{\alpha}-c_2V^{\beta}$, with $0<\alpha<1<\beta$.
The standard fixed-time integral bound gives a settling time no larger than
$T_o$ after the prescribed time transformation. Strong invariance of the
sliding set then gives $\varepsilon_1=\varepsilon_2=0$ for all subsequent
time. Hence $z_2=\Delta_a$ and $\hat d=d$ after $T_o$.
\end{proof}

\subsection{Dynamic compensation and error linearization}
Use
\begin{equation}
 \tau=M_e(q_m)v+C_e(q_m,\dot q_m)\dot q_m-\hat d,
 \qquad v=\ddot q_d+\nu.
 \label{eq:computed_torque}
\end{equation}
Under exact model matching and Assumption~\ref{as:observer_interface},
$\hat d=d$ for $t\geq T_o$, and~\eqref{eq:joint_model_relabelled} gives
$\ddot q_m=v$. Consequently, with
$e=q_m-q_d$ and $z=[e^\top,\dot e^\top]^\top$,
\begin{equation}
 \dot z=Az+B\nu,
 \qquad A=\begin{bmatrix}0&I_n\\0&0\end{bmatrix},
 \qquad B=\begin{bmatrix}0\\I_n\end{bmatrix}.
 \label{eq:linear_error}
\end{equation}
The pair $(A,B)$ is controllable since
$\operatorname{rank}[B,AB]=2n$.

\subsection{Smooth periodic delayed feedback}\label{sec:pdf_design}
Choose $K_0$ so that $A_c=A+BK_0$ is Hurwitz. Let $h>0$ and let
$R_h:[0,2h]\to\mathbb R$ be smooth, nonnegative, zero on $[0,h]$, and
positive on a subset of positive measure in $(h,2h)$. Define
\begin{equation}
 W_c=\int_h^{2h}e^{-A_cs}B R_h(s)B^\top e^{-A_c^\top s}\,ds.
 \label{eq:wc_matrix}
\end{equation}
\begin{lemma}[Positivity of the PDF Gramian]\label{lem:gramian}
If $(A_c,B)$ is controllable, then $W_c\succ0$.
\end{lemma}
\begin{proof}
For any $x\in\mathbb R^{2n}$,
\[
 x^\top W_cx=\int_h^{2h}R_h(s)\|B^\top e^{-A_c^\top s}x\|^2ds\geq0.
\]
If this quantity is zero, continuity implies
$B^\top e^{-A_c^\top s}x=0$ on an interval contained in the set where
$R_h>0$. The left-hand side is analytic in $s$, hence it vanishes for all
$s$. Differentiating at any point gives
$B^\top(A_c^\top)^kx=0$ for every $k\geq0$. Controllability of
$(A_c,B)$ implies $x=0$. Therefore $W_c$ is positive definite.
\end{proof}

For local PDF time $\vartheta=t-T_o$, define the $2h$-periodic gain
\begin{equation}
 K_c(\vartheta)=\begin{cases}
 0,&0\leq\vartheta<h,\\
 R_h(\vartheta)B^\top e^{-A_c^\top\vartheta}W_c^{-1}
 e^{A_c(h-\vartheta)},&h\leq\vartheta<2h,
 \end{cases}
 \label{eq:kc_pdf}
\end{equation}
and extend it periodically. The control input is
\begin{equation}
 \nu(t)=K_0z(t)-K_c(t-T_o)z(t-h),\qquad t\geq T_o.
 \label{eq:pdf_error_input}
\end{equation}
Thus the complete controller is
\begin{equation}
 \tau=M_e\left[\ddot q_d+K_0z-K_c(t-T_o)z(t-h)\right]
 +C_e\dot q_m-\hat d.
 \label{eq:final_controller}
\end{equation}
